\section{Kết quả thực nghiệm và Tự đánh giá}

\subsection{Kết quả kiểm thử trên các bộ dữ liệu}
Tóm tắt kết quả chạy thử nghiệm chương trình trên các trường hợp kiểm thử (test cases):

\begin{table}[h!]
    \centering
    \caption{Bảng kết quả kiểm thử chương trình}
    \label{tab:test_results}
    \begin{tabularx}{\textwidth}{c X c c}
        \toprule
        \textbf{Test Case} & \textbf{Mô tả dữ liệu thử nghiệm} & \textbf{Kết quả mong đợi} & \textbf{Đánh giá} \\
        \midrule
        01 & Dữ liệu mảng rỗng $A = []$ & Trả về $-1$ & Đạt \\
        02 & Phần tử nằm ở đầu mảng & Trả về vị trí $0$ & Đạt \\
        03 & Phần tử không tồn tại trong mảng & Trả về $-1$ & Đạt \\
        04 & Mảng lớn có $10^6$ phần tử & Thời gian $< 0.1\text{s}$ & Đạt \\
        \bottomrule
    \end{tabularx}
\end{table}

\subsection{Bảng tự đánh giá mức độ hoàn thành}
Sinh viên tự đánh giá tỷ lệ hoàn thành các yêu cầu của bài tập:

\begin{table}[h!]
    \centering
    \caption{Tự đánh giá mức độ hoàn thành yêu cầu}
    \label{tab:self_eval}
    \begin{tabularx}{\textwidth}{c X c}
        \toprule
        \textbf{STT} & \textbf{Yêu cầu của đề bài} & \textbf{Mức độ hoàn thành} \\
        \midrule
        1 & Cài đặt đúng thuật toán và cấu trúc dữ liệu yêu cầu & 100\% \\
        2 & Xử lý đầy đủ các trường hợp ngoại lệ (Edge cases) & 100\% \\
        3 & Tối ưu hóa thời gian thực thi và bộ nhớ sử dụng & 100\% \\
        4 & Viết báo cáo rõ ràng, mạch lạc đúng chuẩn quy định & 100\% \\
        \bottomrule
    \end{tabularx}
\end{table}

\subsection{Kết luận}
Chương trình đã hoàn thành đầy đủ các yêu cầu đề ra, chạy ổn định và vượt qua toàn bộ các bộ dữ liệu kiểm thử.
